\section{Player observations and small test domains}\label{sec:perspectives}
This section uses $\mathbf A_I$, with winner constants and owned choices
but no internal role switches. The full-signature comparison is
\cref{view:thm:sym}; \cref{gen:cor:views} shows why these perspectives
also apply intrinsically to all idempotent definable operations.

\subsection{Exactly what one player can observe}

Fix $k\in I$. An observer asks whether $k$ is a possible winner after a
term context has used the input. Put
\[
 S\equiv_k T\quad\Longleftrightarrow\quad
 \forall f\in\operatorname{Pol}_1(\A),\quad
       k\in f(S)\ \Longleftrightarrow\ k\in f(T).
\]
Here the polynomial may repeat the input arbitrarily and use all ground
constants. This is the standard contextual, or Leibniz, congruence of the
logical matrix whose designated values are the sets containing $k$
\cite{BlokPigozzi1989,Font1993}. Computing it is the specific problem.

\begin{theorem}[Exact one-player quotient]\label{view:thm:view}
For nonempty $S,T\subseteq I$,
\begin{equation}\label{view:eq:view}
 S\equiv_k T\quad\Longleftrightarrow\quad
 (k\in S\Longleftrightarrow k\in T)
 \ \land\ \bigwedge_{i\ne k}(S=\{i\}\Longleftrightarrow T=\{i\}).
\end{equation}
There are exactly $K+1$ classes:
\begin{enumerate}
\item $\pv$: every set containing $k$, including $\{k\}$;
\item $a_i$ for each $i\ne k$: the single value $\{i\}$;
\item $\nv$: every ambiguous set omitting $k$.
\end{enumerate}
This is the coarsest congruence preserving the possible-win observation
for $k$.
\end{theorem}
\begin{proof}
For nonempty $S,T$, direct calculation gives
\[
 S\ow r T=\{i\}\ \Longleftrightarrow\
 \begin{cases}
 S=\{i\}\text{ or }T=\{i\},&r=i,\\
 S=T=\{i\},&r\ne i.
 \end{cases}
\]
Also, membership of the owner $r$ is simply the disjunction of its
membership in the two inputs. For $k\ne r$, membership of $k$ is
\[
 (k\in S\text{ or }k\in T)\ \land\ S\ne\{r\}\ \land\ T\ne\{r\}.
\]
These equations show that the right-hand side of~\eqref{view:eq:view} is a
congruence, hence is preserved by every context.

Conversely, the identity context tests membership of $k$. For $i\ne k$,
the context $x\ow i\{k\}$ excludes $k$ exactly when $x=\{i\}$. These
contexts distinguish all the listed cases. The class $\nv$ is inhabited
because there are two distinct roles other than $k$.
\end{proof}

\texttt{sameView\_own} and \texttt{contextualAt\_iff} prove the
congruence and its contextual characterization generically in Lean.
Ground constants suffice to represent arbitrary polynomial parameters
because union realizes every nonempty subset.

Write $C_k=\A/{\equiv_k}$. Its operations have a short explicit form.
For owner $r\ne k$, an input $a_r$ absorbs. Without that absorber, an
input $\pv$ gives $\pv$; two equal atoms $a_i$ give $a_i$; all other
cases give $\nv$. For owner $k$, omit the absorbing-$a_r$ case.
Equivalently, each operation is the join of a simple order: distinct
non-owner atoms join to $\nv$, then comes $\pv$, and an owner atom
$a_r$, when present, sits above everything.

\begin{corollary}[Faithful family of small perspectives]\label{view:cor:product}
The canonical map
\[
 \A\longrightarrow\prod_{k\in I} C_k,\qquad
 S\longmapsto([S]_{\equiv_k})_{k\in I}
\]
is an injective homomorphism and is onto each factor. Thus the
$2^K-1$-valued algebra has a subdirect representation by $K$ factors,
each of size $K+1$.
\end{corollary}
\begin{proof}
Each coordinate is a quotient homomorphism. If two inputs differ, some
$k$ belongs to exactly one of them, and their $k$ coordinates differ.
\end{proof}

This retains possible SPE winners. It is not the Boolean guarantee
representation: $\{k\}$ and every ambiguous set containing $k$ are
\emph{identified} in $C_k$, whereas individual forceability distinguishes
them. The family of all perspectives jointly recovers every membership
bit. Role renaming transports $C_k$ to $C_{\pi k}$ without privileging
one observer.

It also supplies ordinary finite logical matrices: designate $\pv$ in
$C_k$ and interpret consequence by preservation of that designation.
Generating a truth-table proof system for a finite matrix is standard;
the content here is the explicit smallest contextual matrix.

\subsection{Optimal fixed input tests, at every arity}

Call $D\subseteq\Pow(I)$ a complete input test domain if, for every arity
$n$ and every pair of terms $f,g$,
\[
 f=g\text{ on }\Pow(I)^n\quad\Longleftrightarrow\quad f=g\text{ on }D^n.
\]
The domain is chosen once and must work for all terms and all arities.

\begin{theorem}[Exact characterization]\label{view:thm:tests}
A domain $D$ is complete exactly when:
\begin{enumerate}
\item every singleton $\{i\}$ belongs to $D$;
\item for every $k$, some ambiguous $F\in D$ omits $k$.
\end{enumerate}
The second condition says that the intersection of the ambiguous members
of $D$ is empty, with the intersection of an empty family taken as $I$.
\end{theorem}
\begin{proof}
\emph{Sufficiency.} Suppose $f,g$ differ at a valuation $v$, and choose
$k$ belonging to exactly one output. Choose an ambiguous $F\in D$
omitting $k$. Replace each input set $S$ as follows:
\[
 S\longmapsto
 \begin{cases}
 \{k\},&k\in S,\\
 \{i\},&S=\{i\},\ i\ne k,\\
 F,&\text{otherwise}.
 \end{cases}
\]
The replacement belongs to $D$ and has the same $k$-view. All operations
preserve that view, so the outputs still differ at $k$. The replacement
may use the \emph{same} $F$ for every variable. This proof works at any
arity, not just for unary functions.

\emph{Necessity of singleton tests.} The unary classification in \cref{game:thm:unary} realizes all functions
\[
 R(C,E)(S)=\begin{cases}\{i\}\cup E_i,&S=\{i\},\\
 S\cup C,&|S|\geq2,\end{cases}
 \quad E_i\subseteq C\setminus\{i\}.
\]
Fix $i$ and $k\ne i$. Take $C=I$, all rows full except
$E_i=I\setminus\{i,k\}$. This function differs from the constant $I$
only at $\{i\}$. Hence that singleton must be tested.

\emph{Necessity of ambiguity coverage.} For each $k$, the same
classification realizes
\[
 F_k(S)=\begin{cases}S,&|S|=1,\\S\cup\{k\},&|S|\geq2.\end{cases}
\]
This differs from the identity exactly on ambiguous sets omitting $k$.
At least one such input must be tested.
\end{proof}

The sufficiency statement is \texttt{complete\_test\_domain} in Lean.
The necessity statement is \texttt{test\_domain\_necessary} in
\texttt{Synthesis.lean}, using explicit unary synthesis and term
separators. The cardinality corollary below remains a written argument.

\begin{corollary}[Sharp size]\label{view:cor:optimal}
The minimum size of a complete fixed test domain is
\[
 \begin{cases}6,&K=3,\\K+2,&K\geq4.\end{cases}
\]
For three roles, take every singleton and every pair. For four or more,
take every singleton and any two disjoint pairs.
\end{corollary}
\begin{proof}
Every singleton is necessary. One ambiguous set has nonempty
intersection with itself, so at least two are needed. Two disjoint pairs
suffice for $K\geq4$. At $K=3$, every ambiguous set has size at least
two, so any two intersect. All three pairs are necessary and sufficient.
\end{proof}

For a fixed observed player, an even smaller tailored domain consists of
all $K$ singletons and one fixed pair omitting that player. Thus an
equation can be checked in $K$ quotient algebras on $K+1$ states, or on a
single fixed input alphabet of the optimal size above. The latter input
alphabet is not a subalgebra: intermediate outputs may be larger sets.
No claim is made that it is a smaller carrier with the same operations.
